\Needspace{13\baselineskip}
\section{R7: Retrieval without forgetting the dependencies}\label{sec:retrieval}

\Q{R7.Q1}{A useful dependent type abstraction}{How can abstract reachability over universe-polymorphic, dependent signatures remain a sound over-approximation without connecting everything to everything?}

\answer Use a hierarchy of relational type abstractions and an AND-hypergraph of component applications. A declaration with several arguments is not an ordinary edge that becomes usable when just one argument type is reachable. Its application requires all arguments, with shared type parameters and index constraints instantiated consistently. The abstraction can relax these conditions, but must not silently turn a failed approximate match into a proof of concrete impossibility.

At the coarsest useful level, retain the conclusion head, argument heads, binder roles, and whether the result is a proposition, a data type, or a type former. At the next level retain repeated type variables, parameter positions, and a small set of index relationships. For example,
\[
 \mathsf{append}:\mathsf{Vec}\ A\ n\to\mathsf{Vec}\ A\ m
          \to\mathsf{Vec}\ A\ (n+m)
\]
shares the same $A$ across all three vector occurrences, but does not equate $n$ and $m$. Forgetting the arithmetic equation is a legitimate over-approximation; introducing an incorrect equality between the indices is not. Universes should be represented by compatible symbolic constraints or an explicit unknown, never silently specialized to a convenient universe.

A concrete typing derivation must map to an abstract derivation. This is the central soundness obligation of the abstraction. It licenses rejecting an abstractly unreachable expression class only when the abstraction and reachability computation really cover every concrete construction in the claimed grammar. Lambda introduction, local variables, constructors, allowed casts, and all admitted component forms must be included in that coverage. A graph of library constants alone does not cover a language that can synthesize higher-order arguments.

The concrete stage instantiates the declaration's full telescope in a fresh transactional context and checks the resulting arguments and conclusion. Failed unification can identify a spurious abstract connection, but only a suitably justified failure supports a reusable abstraction refinement. A timeout or unsupported elaboration case is not such a justification. Repeatedly rejecting individual concrete candidates, without changing the abstraction, is ordinary concrete filtering rather than type-guided abstraction refinement.

Start with a fast index and a relevance ranking, not a hard global exclusion filter. Fixed top-$k$ retrieval and a path bound of three can be useful speed controls, but change the searched grammar. Include them in the result's bound descriptor. A later fair expansion of the provider set can preserve an enumerative fallback. Output-sensitive behavior matters: no discrimination tree can return a million matches in time independent of the number returned.\cite{proposal,hoogle}

\Q{R7.Q2}{Type classes have already been addressed in TYGAR}{Has type-guided abstraction refinement been combined with type classes, and what should Lean preserve?}

\answer Yes. This is a concrete answer missing from the proposal: Hoogle+ explicitly handles type classes by dictionary passing. Class methods consume dictionary components, while instances supply components that construct dictionaries. Its paper discusses this in its type-class treatment and implementation. The formal setting and supported Haskell subset are nevertheless narrower than dependent Lean; higher-kinded polymorphism is not generally covered. The reported evaluation uses 291 components and 44 queries, so it should not be cited as a demonstrated experiment on several hundred thousand dependent declarations.\cite{hoogle}

The corresponding Lean adaptation should keep instance arguments as typed obligations within the same dependency problem as ordinary arguments. An instance can determine an output type or depend on an unresolved index. Conversely, an ordinary argument may determine the type needed for instance synthesis. Searching instances eagerly for every approximate component match can be expensive; concluding ``no instance exists'' before the relevant metavariables are fixed can be wrong.

Cache instance results under the full class application, universe information, local-instance context, environment fingerprint, and relevant options. Distinguish an established result from a postponed search and from failure under a particular resource bound. Class parameters cannot be reduced to a single \code{(class, type)} pair in general. In particular, output-parameter behavior and multiple class arguments carry information that must remain visible to the elaborator.

A class dictionary is not necessarily a proof-irrelevant object: it may contain operations that affect the program's behavior. Two dictionaries for the same carrier type can implement different comparisons. Observational evaluation and memoization must therefore retain the actual instance identity or a proved equivalence sufficient for the intended use. Class laws should be available as explicit evidence when a pruning rule depends on them.

Finally, ordinary Lean permits weakening and contraction. A component argument need not consume the only available value of its type, and one dictionary can be reused. A Petri-net-inspired representation must encode this ordinary functional behavior rather than accidentally imposing linear-resource discipline. The abstract hypergraph should express possible composition; exact dependent type checking determines whether a particular composition is realizable.

\Q{R7.Q3}{Premise selection for data programs}{Can a premise selector accept a data goal with examples, and can a training corpus be obtained by erasing library definition bodies?}

\answer A wrapper can submit a data goal or an encoding of its signature to an existing relevance model, but that does not establish that the model understands the examples or ranks executable components well. Proof premise selection is trained and evaluated on proof dependencies; program synthesis needs behavior-sensitive selection of functions, constructors, and helpers. LeanHammer is a useful interface and evaluation comparison, not evidence that this transfer already works for Leant2.\cite{leanhammer}

Define a retrieval query containing the exact target type, local telescope, symbols appearing in the contract, and a typed summary of the observations. For lists this can include input/output length relationships, preserved elements, ordering patterns, and whether examples distinguish empty and nonempty cases. These summaries are heuristic features unless justified by a checked theorem. The retrieval reply should be a bounded ordered list of declaration names and optional instantiation sketches, never an accepted proof or a reason to discard all other providers.

A Lean-native baseline should combine a discrimination-tree lookup, symbol overlap, type-shape matching, and a small demand-directed composition search. It needs no external model or server. Compare learned retrieval against this baseline at equal downstream synthesis budgets, rather than comparing retrieval scores alone. The important metric is additional fully certified programs per unit of total work, including the cost of loading the index and evaluating the proposed components.

Erased-definition tasks can supply supervision, but careless construction makes them trivial. Hiding the body text while leaving the original declaration callable allows direct retrieval of the answer. Aliases, generated equation lemmas, implementation rewrites, and closely related definitions can also reveal it. Construct a controlled provider environment that excludes the target and designated revealing dependencies, document the removal policy, and keep a separate unrestricted-library track for realistic reuse. A theorem about a hidden target must not silently reintroduce that target as a callable constant.

Split training and evaluation by definition family and preferably by module or repository, not merely by randomly choosing individual declarations. Near-duplicate variants of the same algorithm should stay in one split. Record which reference-body constants provided supervision, but do not expose those labels at inference time. Likewise, example generation must be separate from an inaccessible reference implementation during synthesis. Universal specification lemmas can serve as contracts only when their statements remain well defined in the controlled environment.

The natural corpus has at least two tracks: construction from a restricted primitive library, which measures synthesis, and reuse from a broad library, which measures retrieval and composition. A third, indexed-family track tests whether type abstraction actually preserves the constraints that matter. No accuracy number for these tracks can be inferred from a proof-premise benchmark; they need direct experiments.
